\documentclass[10pt,a4paper]{article}
\usepackage[utf8]{inputenc}
\usepackage[vietnamese]{babel}
\usepackage{amsmath,amssymb,amsfonts}
\usepackage{geometry}
\geometry{a4paper, left=10mm, right=10mm, top=10mm, bottom=10mm}
\usepackage{tikz}
\usetikzlibrary{shapes,arrows.meta,positioning,calc}
\usepackage[most]{tcolorbox}
\usepackage{array,booktabs,multicol}
\usepackage{fontawesome}
\usepackage{helvet}
\renewcommand{\familydefault}{\sfdefault}
\usepackage{parskip}

\definecolor{primary}{RGB}{112, 48, 160}
\definecolor{teal}{RGB}{0, 128, 128}
\definecolor{orange}{RGB}{230, 108, 44}
\definecolor{green}{RGB}{46, 139, 87}
\definecolor{rose}{RGB}{192, 0, 0}

\pagestyle{empty}

\begin{document}

% --- HEADER BANNER ---
\begin{tcolorbox}[
  enhanced,
  colback=violet!8!white,
  colframe=violet!80!black,
  arc=2.5mm, boxrule=1pt,
  left=3mm, right=3mm, top=2mm, bottom=2mm
]
  \begin{tabular*}{\textwidth}{@{\extracolsep{\fill}} lr}
    {\small\bfseries\color{violet!85!black}\faBook\ SỔ TAY KIẾN THỨC TOÁN 2} & 
    {\small\bfseries\color{teal!85!black}\faGraduationCap\ CHƯƠNG TRÌNH GDPT 2018}
  \end{tabular*}
  \vspace{-1mm}
  \begin{center}
    \rule{0.98\textwidth}{0.6pt}
  \end{center}
  \vspace{-2mm}
  \centering
  {\LARGE\bfseries\color{violet!90!black}\faBookmark\ BÀI 75: ÔN TẬP CHUNG}\\[1mm]
  {\itshape\footnotesize\color{black!75} Đại tổng kết toàn bộ chương trình Toán lớp 2: Số học, Phép tính, Hình học, Đo lường}
\end{tcolorbox}

\vspace{-2mm}

\noindent
\begin{minipage}[t]{0.49\textwidth}\vspace{0pt}
\begin{tcolorbox}[
  enhanced,
  colback=violet!4!white,
  colframe=violet!75!black,
  arc=2mm, boxrule=0.8pt,
  left=2mm, right=2mm, top=1.8mm, bottom=1.8mm,
  title={\footnotesize\bfseries\color{white}\faKey\ 1. Trụ cột 1: Số học và Phép tính}
]
{\footnotesize
\begin{itemize}\setlength{\itemsep}{1pt}\setlength{\leftskip}{-2.5mm}
  \item \textbf{Phạm vi số:} Các số đến $100$ (HK1) và đến $1\,000$ (HK2).
  \item Cấu tạo số $\overline{abc} = a00 + b0 + c$; so sánh và thứ tự.
  \item \textbf{4 phép tính nền tảng:}
  \begin{itemize}\setlength{\itemsep}{0.5pt}\setlength{\leftskip}{-2mm}
    \item Phép cộng, phép trừ có nhớ (không quá 1 lần) trong phạm vi $100$ và $1000$.
    \item Phép nhân, phép chia: Thuộc lòng bảng nhân 2, 5 và bảng chia 2, 5.
  \end{itemize}
  \item Tính giá trị biểu thức; tìm thành phần chưa biết ($x$).
\end{itemize}
}
\end{tcolorbox}
\end{minipage}
\hfill
\begin{minipage}[t]{0.49\textwidth}\vspace{0pt}
\begin{tcolorbox}[
  enhanced,
  colback=teal!4!white,
  colframe=teal!75!black,
  arc=2mm, boxrule=0.8pt,
  left=2mm, right=2mm, top=1.8mm, bottom=1.8mm,
  title={\footnotesize\bfseries\color{white}\faSliders\ 2. Trụ cột 2: Hình học và Đo lường}
]
{\footnotesize
\begin{itemize}\setlength{\itemsep}{1pt}\setlength{\leftskip}{-2.5mm}
  \item \textbf{Hình học:}
  \begin{itemize}\setlength{\itemsep}{0.5pt}\setlength{\leftskip}{-2mm}
    \item Điểm, đoạn thẳng, đường thẳng, 3 điểm thẳng hàng.
    \item Đường gấp khúc và tính độ dài đường gấp khúc.
    \item Hình phẳng (tam giác, tứ giác) và hình khối (hộp chữ nhật, lập phương, trụ, cầu).
  \end{itemize}
  \item \textbf{Đại lượng đo lường:}
  \begin{itemize}\setlength{\itemsep}{0.5pt}\setlength{\leftskip}{-2mm}
    \item $cm, dm, m, km$ ($1m = 10dm = 100cm; 1km = 1\,000m$).
    \item $kg$ (cân nặng), $l$ (dung tích), giờ/phút/ngày/tháng.
    \item Tiền Việt Nam ($100 \dots 1\,000$ đồng).
  \end{itemize}
\end{itemize}
}
\end{tcolorbox}
\end{minipage}

\vspace{-1.5mm}
\begin{tcolorbox}[
  enhanced,
  colback=orange!4!white,
  colframe=orange!80!black,
  arc=2mm, boxrule=0.8pt,
  left=2mm, right=2mm, top=1.5mm, bottom=1.5mm,
  title={\footnotesize\bfseries\color{white}\faSitemap\ 3. Bản đồ tư duy kim cương: Toàn cảnh chương trình Toán lớp 2}
]
\centering
\begin{tikzpicture}[x=1.65cm, y=0.75cm]
  \tikzset{arr/.style={-{Stealth[scale=0.8]}}, every node/.append style={align=center}}
  \node[rectangle, rounded corners=2pt, draw=violet!90!black, fill=violet!25, font=\footnotesize\bfseries] (root) at (0,0) {TOÁN LỚP 2\\GDPT 2018};
  \node[rectangle, rounded corners=2pt, draw=teal!80!black, fill=teal!15, font=\scriptsize\bfseries] (p1) at (2.4,0.9) {1. SỐ HỌC VÀ TÍNH TOÁN\\Cộng, Trừ, Nhân (2,5), Chia (2,5)};
  \node[rectangle, rounded corners=2pt, draw=orange!80!black, fill=orange!15, font=\scriptsize\bfseries] (p2) at (2.4,0.3) {2. HÌNH HỌC VÀ ĐO LƯỜNG\\Đường gấp khúc, Khối trụ/cầu, $m, km$};
  \node[rectangle, rounded corners=2pt, draw=green!70!black, fill=green!15, font=\scriptsize\bfseries] (p3) at (2.4,-0.3) {3. THỐNG KÊ VÀ XÁC SUẤT\\Kiểm đếm, Biểu đồ tranh, Khả năng};
  \node[rectangle, rounded corners=2pt, draw=red!75!black, fill=red!15, font=\scriptsize\bfseries] (p4) at (2.4,-0.9) {4. GIẢI TOÁN CÓ LỜI VĂN\\Nhiều hơn, ít hơn, chia đều, gấp lần};

  \draw[arr, thick] (root) |- (p1);
  \draw[arr, thick] (root) |- (p2);
  \draw[arr, thick] (root) |- (p3);
  \draw[arr, thick] (root) |- (p4);
\end{tikzpicture}
\end{tcolorbox}

\vspace{-1.5mm}
\noindent
\begin{minipage}[t]{0.49\textwidth}\vspace{0pt}
\begin{tcolorbox}[
  enhanced,
  colback=green!4!white,
  colframe=green!60!black,
  arc=2mm, boxrule=0.8pt,
  left=2mm, right=2mm, top=1.8mm, bottom=1.8mm,
  title={\footnotesize\bfseries\color{white}\faCheckCircle\ 4. Bài tập tổng hợp cuối năm mẫu}
]
{\footnotesize
\textbf{1. Tính:}\\
- $452 + 286 = \mathbf{738}$; \quad $745 - 328 = \mathbf{417}$.\\
- $5 \times 8 + 60 = 40 + 60 = \mathbf{100}$.\\
- $18 : 2 + 35 = 9 + 35 = \mathbf{44}$.\\
\textbf{2. Bài toán lời văn tổng hợp:}\\
Một cửa hàng buổi sáng bán được $235\,kg$ gạo, buổi chiều bán được nhiều hơn buổi sáng $45\,kg$ gạo. Hỏi cả hai buổi cửa hàng bán được bao nhiêu ki-lô-gam gạo?\\
\textit{Giải:} Buổi chiều bán được: $235 + 45 = 280\,(kg)$.\\
Cả hai buổi bán được: $235 + 280 = \mathbf{515\,(kg)}$.
}
\end{tcolorbox}
\end{minipage}
\hfill
\begin{minipage}[t]{0.49\textwidth}\vspace{0pt}
\begin{tcolorbox}[
  enhanced,
  colback=rose!4!white,
  colframe=rose!70!black,
  arc=2mm, boxrule=0.8pt,
  left=2mm, right=2mm, top=1.8mm, bottom=1.8mm,
  title={\footnotesize\bfseries\color{white}\faGraduationCap\ 5. Hành trang vững bước vào Lớp 3}
]
{\footnotesize
\begin{itemize}\setlength{\itemsep}{1.5pt}\setlength{\leftskip}{-2.5mm}
  \item \textbf{Tự tin làm chủ kiến thức:} Đọc viết số đến $1000$, thành thạo cộng trừ có nhớ và thuộc lòng bảng nhân/chia 2 và 5.
  \item \textbf{Kĩ năng làm bài thi:} Trình bày sạch đẹp, đặt tính thẳng cột, kiểm tra lại đáp số và ghi đầy đủ tên đơn vị.
  \item Chúc các em học sinh hoàn thành xuất sắc chương trình Toán lớp 2 và sẵn sàng bước vào Toán 3!
\end{itemize}
}
\end{tcolorbox}
\end{minipage}


\end{document}
